% Generated draft: requires mathematical and source review.
\documentclass[11pt]{article}
\usepackage[margin=25mm]{geometry}
\usepackage{fontspec}
\setmainfont{Latin Modern Roman}
\usepackage{amsmath,amssymb,mathtools,unicode-math}
\setmathfont{Latin Modern Math}
\usepackage{enumitem,needspace,xurl,hyperref}
\setlength{\parindent}{0pt}
\setlength{\parskip}{5pt}
\setlength{\emergencystretch}{3em}
\newcommand{\N}{\mathbb{N}}
\newcommand{\Z}{\mathbb{Z}}
\newcommand{\Q}{\mathbb{Q}}
\newcommand{\R}{\mathbb{R}}
\newcommand{\C}{\mathbb{C}}
\newcommand{\F}{\mathbb{F}}
\newcommand{\A}{\mathbb{A}}
\newcommand{\PP}{\mathbb{P}}
\newcommand{\E}{\mathbb{E}}
\newcommand{\eps}{\varepsilon}
\newcommand{\dd}{\,\mathrm{d}}
\begin{document}
\section*{UnsolvedMath ID: 30003729}
\textbf{Source problem code:} catalogue entry
\subsection{Definitions and mathematical statement}
Let \(\eta\) denote the perturbation parameter in the source's randomly perturbed periodic coercive Hamilton--Jacobi setting. Write \(\overline H\) for the unperturbed ergodic constant and \(\overline H^\eta\) for the perturbed ergodic constant. These are the quantities designated by the supplied record; their defining equations and existence hypotheses are not supplied. The original statement writes these quantities as \(H\) and \(H^\eta\); the notation here follows the cleaned statement.

For nonzero admissible values of \(\eta\), define the normalized change by
\[
D(\eta):=\frac{\overline H^\eta-\overline H}{\eta}.
\]
The source asks for the behavior of \(\overline H^\eta\) in general settings, especially the asymptotic behavior of \(D(\eta)\) as \(\eta\to0\). No existence or value of a limit is asserted.

This is an unresolved formulation of a research target rather than a fully specified quantified conjecture. The record does not identify the Hamiltonian, spatial dimension or domain, equation, boundary conditions, random perturbation mechanism or probability law, or the meaning and existence conditions of the ergodic constants. It also does not specify the admissible parameter range, the direction of approach to zero, any dependence on a fixed momentum, or the meaning of convergence if the constants remain random. Periodicity and coercivity are stated, but their precise hypotheses and their applicability before and after perturbation are not given. Consequently, a precise class of equations and a universal assertion about this quotient cannot be recovered without further source information.
\subsection{Short English statement}
The source asks how random perturbations change the ergodic constant of a periodic coercive Hamilton--Jacobi equation, particularly the change divided by \(\eta\) as \(\eta\to0\). The equation and perturbation model are unspecified, so this draft preserves the research target and records the missing hypotheses. The supplied partially solved label is a dated curation assessment, not an independently verified conclusion.
\subsection{Sources}
\begin{itemize}
\item Supplied local UnsolvedMath registry record, permanent ID \(30003729\), title: \emph{Perturbed Ergodic Constants in Hamilton--Jacobi Equations}; collection: Oberwolfach Reports Open Problems; exact source problem code: \texttt{catalogue entry}. The record gives proposed year 2018 and names no proposer.
\item Supplied source citation: \emph{Interplay of Analysis and Probability in Applied Mathematics} (2019), Oberwolfach Reports, pp. 311--381, DOI 10.4171/owr/2018/6. Supplied source URL: \url{https://doi.org/10.4171/owr/2018/6}. The background describes Oberwolfach Reports as documenting mathematical workshops and their open-problem discussions. No problem-specific page is supplied.
\item Supplied literature source, labeled ``original OWR report'': \url{https://doi.org/10.4171/OWR/2018/6}.
\item Supplied literature source, labeled ``Cardaliaguet--Le Bris--Souganidis, perturbation problem'': \url{https://doi.org/10.1016/j.matpur.2018.03.005}.
\item Literature assessment dated exactly \texttt{2026-08-21}: the supplied curation labels the problem partially solved and reports perturbation expansions and periodic-approximation results in important Hamilton--Jacobi settings, while describing broader behavior as unresolved. This is supplied dated curation, not independently verified literature assessment. The background explicitly requests independent verification of time-sensitive status and cited claims.
\item The registry labels the statement \texttt{corrected\_verified} and reports a manual semantic, typographic, and LaTeX review claiming that the cleaned statement is self-contained and preserves the target. This is supplied review metadata; it does not resolve the missing mathematical hypotheses identified here.
\end{itemize}
\end{document}
